% BHU PYQ -- Psychology: Foundation of Psychology (2024-25 NEP)
\documentclass[11pt,a4paper]{article}
\usepackage[a4paper,top=2.5cm,bottom=2.5cm,left=2.8cm,right=2.8cm,headheight=14pt]{geometry}
\usepackage{amsmath,amssymb}
\usepackage{fontspec}
\setmainfont{Kohinoor Devanagari}
\usepackage{microtype,fancyhdr,enumitem,hyperref,lastpage,parskip}

\hypersetup{colorlinks=true,urlcolor=black,linkcolor=black,pdftitle={PSYMJ11: Foundation of Psychology -- BHU NEP 2024-25}}

\pagestyle{fancy}\fancyhf{}
\renewcommand{\headrulewidth}{0.4pt}\renewcommand{\footrulewidth}{0.4pt}
\fancyhead[L]{\small \textit{Banaras Hindu University}}
\fancyhead[R]{\small \textit{PSYMJ11: Foundation of Psychology (2024-25)}}
\fancyfoot[C]{\small Page \thepage\ of \pageref{LastPage}}

\newlist{parts}{enumerate}{2}
\setlist[parts,1]{label=(\alph*),leftmargin=*,itemsep=4pt,topsep=3pt}
\setlist[parts,2]{label=(\roman*),leftmargin=*,itemsep=2pt,topsep=2pt}

\newcommand{\pts}[1]{\hfill \textit{[#1]}}

\begin{document}

\begin{center}
  {\large\bfseries BANARAS HINDU UNIVERSITY}\\[2pt]
  {\normalsize B. A. (Hons.) Semester I Examination, 2024-25 (NEP)}\\[6pt]
  \rule{0.8\linewidth}{0.5pt}\\[6pt]
  {\large\bfseries PSYCHOLOGY}\\[4pt]
  {\large\bfseries Paper Code: PSYMJ11 : Foundation of Psychology}\\[4pt]
  {\normalsize Time: 2:30 Hours \hfill Full Marks: 70}\\[2pt]
  {\small REF NO. CE/Conf./D/223}
\end{center}
\rule{\linewidth}{0.4pt}
\smallskip

\noindent \textit{Instructions: Answer all questions. Write your Roll No.\ at the top immediately on receipt of this question paper.}

\smallskip
\noindent \textit{सभी प्रश्नों के उत्तर दीजिए।}

\bigskip

\section*{SECTION -- A / खण्ड -- अ}
\noindent \textit{Note: Answer the following questions very briefly. Each question carries 2 marks.}\pts{5 $\times$ 2 = 10}

\noindent \textit{निम्नलिखित प्रश्नों के अति संक्षिप्त उत्तर दीजिए। प्रत्येक प्रश्न 2 अंकों का है।}

\begin{enumerate}[label=\textbf{\arabic*.},leftmargin=*,itemsep=8pt]
  \item 
  \begin{parts}
    \item What do you understand by the concept of acculturation?\pts{2}
    
    उत्संसकृतीकरण (परसंस्कृतिग्रहण) की अवधारणा से आप क्या समझते हैं?

    \item Define the term Illusion.\pts{2}
    
    भ्रम (Illusion) को परिभाषित कीजिए।

    \item What do you understand by the concept of Experiment?\pts{2}
    
    प्रयोग की अवधारणा से आप क्या समझते हैं?

    \item What is the hearing range of human ear?\pts{2}
    
    मानव कान की श्रवण परास (Hearing Range) कितनी होती है?

    \item What do you understand by Figure and Background law of perception?\pts{2}
    
    प्रत्यक्षीकरण के आकृति एवं पृष्ठभूमि नियम से आप क्या समझते हैं?
  \end{parts}
\end{enumerate}

\bigskip

\section*{SECTION -- B / खण्ड -- ब}
\noindent \textit{Note: Give short answer to each of the following questions. Each question carries 10 marks.}\pts{3 $\times$ 10 = 30}

\noindent \textit{निम्नलिखित प्रत्येक प्रश्न का संक्षिप्त उत्तर दीजिए। प्रत्येक प्रश्न 10 अंकों का है।}

\begin{enumerate}[label=\textbf{\arabic*.},leftmargin=*,start=2,itemsep=12pt]

  \item Explain the basic structure of the human ear.\pts{10}
  
  मानव कान की मौलिक संरचना की व्याख्या कीजिए।

  \begin{center}\textbf{OR / अथवा}\end{center}

  Write a short note on the subfields of Psychology.\pts{10}
  
  मनोविज्ञान के उप-क्षेत्रों पर एक संक्षिप्त टिप्पणी लिखिए।

  \item Explain the major perspectives of modern psychology.\pts{10}
  
  आधुनिक मनोविज्ञान के प्रमुख दृष्टिकोणों की व्याख्या कीजिए।

  \begin{center}\textbf{OR / अथवा}\end{center}

  Define Psychology and discuss the major goals of psychology.\pts{10}
  
  मनोविज्ञान को परिभाषित कीजिए तथा मनोविज्ञान के प्रमुख लक्ष्यों की चर्चा कीजिए।

  \item What do you understand by Attention? Explain the various types of attention.\pts{10}
  
  अवधान (Attention) से आप क्या समझते हैं? अवधान के विभिन्न प्रकारों की व्याख्या कीजिए।

  \begin{center}\textbf{OR / अथवा}\end{center}

  Discuss the scientific nature of psychology as a subject.\pts{10}
  
  एक विषय के रूप में मनोविज्ञान की वैज्ञानिक प्रकृति की चर्चा कीजिए।

\end{enumerate}

\bigskip

\section*{SECTION -- C / खण्ड -- स}
\noindent \textit{Note: Answer each of the following questions in detail. Each question carries 15 marks.}\pts{2 $\times$ 15 = 30}

\noindent \textit{निम्नलिखित प्रत्येक प्रश्न का विस्तृत उत्तर दीजिए। प्रत्येक प्रश्न 15 अंकों का है।}

\begin{enumerate}[label=\textbf{\arabic*.},leftmargin=*,start=5,itemsep=14pt]

  \item What do you understand by Observation? Discuss the different types of observation in detail.\pts{15}
  
  प्रेक्षण (Observation) से आप क्या समझते हैं? प्रेक्षण के विभिन्न प्रकारों की विस्तृत चर्चा कीजिए।

  \begin{center}\textbf{OR / अथवा}\end{center}

  Write a note on the socio-cultural basis of human behaviour.\pts{15}
  
  मानव व्यवहार के सामाजिक-सांस्कृतिक आधारों पर एक टिप्पणी लिखिए।

  \item Define the concept of Perception and discuss the various laws of perceptual organization.\pts{15}
  
  प्रत्यक्षीकरण की अवधारणा को परिभाषित कीजिए तथा प्रत्यक्षीकरणात्मक संगठन के विभिन्न नियमों की चर्चा कीजिए।

  \begin{center}\textbf{OR / अथवा}\end{center}

  Define the concept of Sensation and explain the basic structure of the human eye.\pts{15}
  
  संवेदना की अवधारणा को परिभाषित कीजिए तथा मानव आँख की मौलिक संरचना की व्याख्या कीजिए।

\end{enumerate}

\end{document}
